\documentclass[11pt,a4paper]{article}

\usepackage[utf8]{inputenc}
\usepackage[T1]{fontenc}
\usepackage[french]{babel}
\usepackage[a4paper,margin=2cm]{geometry}
\usepackage{amsmath,amssymb,amsthm}
\usepackage{lmodern}
\usepackage{xcolor}

\title{Corrigé détaillé — DS Algèbre linéaire}
\date{}

\newcommand{\commentaire}[1]{\textcolor{blue}{\textit{#1}}}

\begin{document}

\maketitle

%==================================================
\section*{Problème 1}

\subsection*{Rappel de cours}

Un sous-espace vectoriel est un ensemble stable par combinaison linéaire.  
Une somme est directe si l’intersection est réduite à $\{0\}$.

%==================================================
\subsection*{Partie A}

\textbf{1. Montrer que $E_n$ est un sous-espace vectoriel}

Soient $x,y \in E_n$ et $\lambda \in \mathbb{R}$.

Par définition :
\[
x=(x_1,\dots,x_n,0), \quad y=(y_1,\dots,y_n,0)
\]

Alors :
\[
x+y=(x_1+y_1,\dots,x_n+y_n,0)
\]
et :
\[
\lambda x=(\lambda x_1,\dots,\lambda x_n,0)
\]

Donc $x+y \in E_n$ et $\lambda x \in E_n$.

Ainsi, $E_n$ est stable par combinaison linéaire.

\[
\boxed{E_n \text{ est un sous-espace vectoriel}}
\]

\commentaire{Beaucoup d’élèves oublient d’écrire explicitement la stabilité complète.}

%--------------------------------------------------

\textbf{2. Dimension}

La famille $(e_1,\dots,e_n)$ est libre et génératrice de $E_n$.

Donc :
\[
\boxed{\dim(E_n)=n}
\]

\commentaire{Réponse rapide attendue, inutile de surdévelopper.}

%--------------------------------------------------

\textbf{3. Codimension}

Par définition :
\[
\mathrm{codim}(E_n)=\dim(E_{n+1})-\dim(E_n)
\]

Or :
\[
\dim(E_{n+1})=n+1, \quad \dim(E_n)=n
\]

Donc :
\[
\boxed{\mathrm{codim}(E_n)=1}
\]

\commentaire{Attention aux oublis de justification du calcul de dimension.}

%==================================================
\subsection*{Partie B}

\textbf{4. Décomposition}

Soit $x \in E_{n+1}$.

On écrit :
\[
x=(x_1,\dots,x_n,x_{n+1})
\]

On pose :
\[
y=(x_1,\dots,x_n,0) \in E_n
\]
\[
z=x_{n+1}e_{n+1} \in \mathrm{Vect}(e_{n+1})
\]

Alors :
\[
x=y+z
\]

De plus, si :
\[
x=y+z=0
\Rightarrow y=-z
\]
alors $y$ et $z$ ont des formes incompatibles sauf si nuls.

Donc :
\[
E_{n+1}=E_n \oplus \mathrm{Vect}(e_{n+1})
\]

\commentaire{Point clé du sujet — valoriser fortement.}

%--------------------------------------------------

\textbf{5. Construction inductive}

Par récurrence :

\[
E_1=\mathrm{Vect}(e_1)
\]

et :
\[
E_{n+1}=E_n \oplus \mathrm{Vect}(e_{n+1})
\]

Donc :
\[
E_n=\mathrm{Vect}(e_1,\dots,e_n)
\]

\commentaire{Certains élèves oublient la récurrence explicite.}

%--------------------------------------------------

\textbf{6. Caractérisation des bases}

Un vecteur hors de $E_n$ a une dernière coordonnée non nulle.

Donc il n’est pas combinaison linéaire des vecteurs de $E_n$.

Ainsi, il complète toute base.

\commentaire{Lien implicite avec théorème de la base incomplète.}

%==================================================
\subsection*{Partie C}

\textbf{7. Croissance}

\[
F_n \subset F_{n+1} \Rightarrow \dim(F_n) \leq \dim(F_{n+1})
\]

\textbf{8. Majoration}

\[
F_n \subset E_n \Rightarrow \dim(F_n) \leq n
\]

\textbf{9. Stabilité}

Si :
\[
\dim(F_{n+1})=\dim(F_n)
\]

et :
\[
F_n \subset F_{n+1}
\]

alors :
\[
\boxed{F_n=F_{n+1}}
\]

\commentaire{Résultat fondamental en dimension finie.}

%==================================================
\subsection*{Partie D}

\textbf{11. Espace vectoriel}

Soient $x,y \in E$.

Ils appartiennent à un certain $E_n$.

Donc $x+y \in E_n \subset E$.

\textbf{12. Non finiment engendré}

Supposons $E$ engendré par une famille finie.

Alors elle est contenue dans un $E_N$.

Donc impossible d’engendrer tout $E$.

\[
\boxed{E \text{ n’est pas de dimension finie}}
\]

\commentaire{Question discriminante.}

%==================================================
\section*{Problème 2}

\subsection*{Rappel}

Une famille est libre si :
\[
\sum \lambda_i u_i =0 \Rightarrow \lambda_i=0
\]

%--------------------------------------------------

\textbf{1.}

Par définition, chaque vecteur est hors du précédent.

Donc aucune dépendance.

\[
\boxed{\text{famille libre}}
\]

%--------------------------------------------------

\textbf{2.}

Contre-exemple :

$(u,v)$ libre mais $(v,u)$ non libérée.

\commentaire{Beaucoup oublient de donner un exemple explicite.}

%--------------------------------------------------

\textbf{4. Réordonnancement}

On ajoute les vecteurs successivement hors du sous-espace engendré.

\commentaire{Méthode algorithmique implicite.}

%--------------------------------------------------

\textbf{7. Caractérisation}

Libérée + génératrice :
\[
\Rightarrow \text{base}
\]

\commentaire{Question centrale du problème.}

%==================================================
\section*{Problème 3}

\subsection*{Partie A}

\[
A_n=u+nv,\quad B_n=v+nu
\]

\[
\overrightarrow{A_nB_n}=B_n-A_n=(n-1)u+(1-n)v
\]

\[
\overrightarrow{A_nC_n}=v_n
\]

Les vecteurs sont non colinéaires car $u,v$ le sont.

\[
\Rightarrow \text{triangle}
\]

\commentaire{Erreur fréquente : inversion des coefficients.}

%--------------------------------------------------

\subsection*{Partie B}

Tous les vecteurs sont combinaisons de $u,v$.

Donc :
\[
\boxed{\text{plan invariant } = \mathrm{Vect}(u,v)}
\]

%--------------------------------------------------

\subsection*{Partie C}

\[
\overrightarrow{A_nB_n}=(n-1)(u-v)
\]

Donc direction constante.

\[
\boxed{\text{toujours parallèle à } u-v}
\]

\commentaire{Question piège classique.}

%--------------------------------------------------

\subsection*{Partie D}

\[
G_n=\frac{1}{3}(2+n)(u+v)
\]

Suite affine en $n$.

Donc déplacement sur une droite.

\commentaire{Lecture géométrique attendue.}

%--------------------------------------------------

\subsection*{Partie E}

\[
S_n=\mathrm{Vect}(u_n,v_n)=\mathrm{Vect}(u,v)
\]

Donc indépendant de $n$.

\[
\boxed{\text{plan invariant}}
\]

\end{document}